\documentclass[a4paper,12pt]{article}

% ---------- Geometry & typography ----------
\usepackage{fullpage}          % standard ~1 inch margins
\usepackage[T1]{fontenc}
\usepackage{mathptmx}          % Times-like serif (roman, maths)
\usepackage[scaled]{helvet}    % Times-compatible sans (headings optional)
\usepackage{courier}           % Typewriter
\usepackage{textcomp}

% ---------- Layout helpers ----------
\usepackage{graphicx}
\graphicspath{{figures/}}
\usepackage{amsmath}
\usepackage{booktabs}
\usepackage{array}
\usepackage{longtable}
\usepackage{caption}
\usepackage{setspace}
\usepackage{titlesec}
\usepackage{hyperref}
\usepackage{placeins} % \FloatBarrier: keep floats from crossing major sections
\usepackage{siunitx}
\usepackage{xurl}
\usepackage{indentfirst}       % ensure first paragraph after every heading is indented

\hypersetup{
    colorlinks=true,
    linkcolor=black,
    citecolor=black,
    urlcolor=blue
}

% ---------- Line spacing: fixed value 20 pt (Word equivalent) ----------
% With 12 pt body, default baseline is about 14.5 pt; 20/14.5 ~= 1.379
\setstretch{1.379}

% ---------- Paragraph indentation: 2 Chinese-width characters ~ 2em ----------
\setlength{\parindent}{2em}
\setlength{\parskip}{0pt}
\emergencystretch=2em

% ---------- Heading font sizes per template ----------
% First level: Times New Roman, 16 pt
% Second level: Times New Roman, 14 pt
% Third level: Times New Roman, 12 pt
% Text part: Times New Roman, 12 pt (provided by documentclass option + mathptmx)
\titleformat{\section}
  {\normalfont\bfseries\fontsize{16pt}{20pt}\selectfont}
  {\thesection}{1em}{}
\titleformat{\subsection}
  {\normalfont\bfseries\fontsize{14pt}{20pt}\selectfont}
  {\thesubsection}{1em}{}
\titleformat{\subsubsection}
  {\normalfont\bfseries\fontsize{12pt}{20pt}\selectfont}
  {\thesubsubsection}{1em}{}

% Compact headings: no extra ``blank line'' before titles; spacing in fixed pt
% (left, before-title, after-title). After-title keeps one logical line gap to body.
\titlespacing*{\section}{0pt}{4pt plus 1pt minus 1pt}{6pt plus 1pt minus 1pt}
\titlespacing*{\subsection}{0pt}{3pt plus 1pt minus 1pt}{5pt plus 1pt minus 1pt}
\titlespacing*{\subsubsection}{0pt}{2pt plus 1pt minus 1pt}{4pt plus 1pt minus 1pt}

% Captions: same size as body (12 pt); only ``Figure N.'' / ``Table N.'' bold (IEEE-style)
\captionsetup{
  font=normalsize,
  labelfont=bf,
  textfont=normalfont,
  labelsep=period,
  skip=6pt
}
% Prefer floats on text pages; avoid dedicated float-only pages where possible
\setcounter{topnumber}{6}
\setcounter{bottomnumber}{6}
\setcounter{totalnumber}{10}
\renewcommand{\topfraction}{0.90}
\renewcommand{\bottomfraction}{0.90}
\renewcommand{\textfraction}{0.07}
\renewcommand{\floatpagefraction}{0.88}
% Tall figures shrink to fit roughly half a text page (width unchanged when height binds)
\newcommand{\ThesisIncludeGraphics}[2][0.82]{%
  \includegraphics[width=#1\linewidth,height=0.45\textheight,keepaspectratio]{#2}%
}

\begin{document}
\hypersetup{pageanchor=false}

\pagenumbering{roman}

\section*{Abstract}
\addcontentsline{toc}{section}{Abstract}

Semantic communication shifts the objective of communication systems from bit-level fidelity to meaning-level fidelity, making it a promising direction for post-Shannon wireless communication. This project investigates a DeepSC-style Transformer-based semantic transceiver for text transmission over AWGN and Rayleigh fading channels. The study focuses on three improvement directions: SNR-aware attention, adaptive semantic-loss weighting, and language-model-based post-refinement. A baseline system is first reproduced, after which several modified variants are evaluated under the same data split, SNR grid, and metric definitions.

The results show that the improved transceiver gives measurable gains under low-SNR AWGN conditions. At \SI{-10}{dB}, BLEU-4 increases from 0.4189 to 0.4329, and Word Accuracy increases from 0.4425 to 0.4773. However, the semantic-loss ablation shows that removing the semantic-loss term improves BLEU-4 across all tested SNR points, indicating that the current semantic regularization is not well aligned with token-level reconstruction. Under Rayleigh fading, the gains are limited and mixed, suggesting that scalar SNR conditioning is insufficient for strongly time-varying channels. Confidence-driven BERT refinement also fails to provide reliable net gains because of over-correction. In comparison, beam search decoding gives the most consistent inference-stage improvement, reaching BLEU-4 = 0.4552 at \SI{-10}{dB} without retraining.

Overall, the project suggests that robust channel-aware design and decoding-strategy optimization are more reliable than unconstrained external post-refinement in this setting.

\vspace{0.5em}
\noindent\textbf{Keywords:} Semantic Communication; DeepSC; SNR-Adaptive Attention;
Adaptive Semantic Loss; BERT Masked Refinement; Beam Search; Rayleigh Fading;
Transformer.
\clearpage
\section*{Acknowledgements}
\addcontentsline{toc}{section}{Acknowledgements}

I would like to express my sincere gratitude to my supervisors for the continuous guidance throughout this project, including topic definition, experiment design, technical analysis, and scientific writing. Their patience with iterative redesign, detailed feedback on early drafts, and insistence on evidence-driven conclusions significantly shaped the final form of this thesis. I am also grateful to members of our research group for valuable discussions on debugging, model evaluation, and interpretation of negative experimental findings. Their practical suggestions helped me avoid several dead ends during the implementation phase and greatly improved the reliability of the reported results.

Special thanks are due to the open-source community and to the developers of the deep-learning and scientific Python ecosystem, whose libraries and frameworks made the full engineering pipeline feasible on a single workstation. The availability of high-quality pretrained language models, tokenizers, and evaluation utilities reduced a considerable amount of engineering overhead and allowed more time to be spent on model design, evaluation, and analysis. I am likewise grateful to the authors of the datasets used in this project, whose careful data curation underpins every experiment reported here.

Finally, I deeply appreciate my family and friends for their encouragement and support during the entire project period. Their emotional support was essential during the most difficult episodes and during the long training runs. This project would not have been possible without the combined contribution of all of the people above.

\clearpage
\tableofcontents
\clearpage
\hypersetup{pageanchor=true}
\pagenumbering{arabic}

% Each major numbered section (1--4) starts on a new page.
\section{Introduction}
\subsection{Project Background and Motivation}

Conventional communication systems are firmly rooted in Shannon's mathematical theory of communication, where the main design objective is accurate bit delivery through noisy channels \cite{shannon1948, shannon1948part2}. This objective has driven decades of successful progress in source coding, channel coding, modulation, and multiple-access schemes. However, an increasing number of emerging intelligent applications, such as human--machine dialog systems, distributed artificial-intelligence (AI) collaboration, autonomous driving control loops, and edge semantic sensing, do not strictly require exact bitwise reconstruction if the underlying meaning of the message remains useful for the downstream task. This conceptual gap between syntactic fidelity and semantic usefulness has motivated the recent resurgence of semantic communication, in which the design objective is redefined from bit-level fidelity to meaning-level fidelity \cite{carnap1954, bao2011semantic}.

Recent research has shown that deep neural networks can learn task-oriented semantic representations and perform end-to-end semantic transmission over realistic channel models \cite{xie2021deepsc, lu2024survey}. In particular, Transformer-based semantic transceivers have become a mainstream design choice thanks to their strong contextual modelling capability, their parallelizable structure, and their compatibility with pretrained language priors \cite{vaswani2017attention}. Despite the promising progress reported in the literature, several deployment-oriented issues remain unsolved. Low-SNR and fading conditions still cause severe semantic degradation; model components that appear theoretically beneficial may become unstable in practice when interacting with channel noise and other components; external post-processing models driven by language priors may introduce additional error rather than effective correction; and the behavior observed under an AWGN benchmark does not necessarily transfer to Rayleigh fading scenarios that are closer to real wireless propagation.

This project is motivated by a very practical engineering question: can one construct a semantic transceiver that is not only theoretically improved, but also experimentally validated through a rigorous iterative process whose outcomes, whether positive or negative, are reproducible and traceable? The emphasis here is not merely on ``adding more neural modules'' on top of an existing pipeline. Instead, this project aims to establish a full hypothesis--diagnosis--redesign--validation loop that identifies both effective design choices and ineffective ones.

From an undergraduate research perspective, this motivation has two very practical implications. First, the project should remain reproducible and implementation-feasible on a single GPU workstation, so that each design decision can be traced back to concrete code changes and to measurable outcomes in saved evaluation files. Second, the final report should not merely highlight improvements but also explicitly document negative findings and their root causes, because engineering decisions in real systems are typically constrained by robustness, interpretability, latency, and memory footprint rather than by isolated best-case numbers.

In this context, the project combines model-side architectural design, training-side objective tuning, and inference-side strategy optimization. This integrated view is necessary because semantic communication quality is jointly affected by representation learning behavior, channel perturbation statistics, and decoding policy interactions. A single-module improvement may appear promising in isolation yet fail after being integrated with other modules or evaluated under a different channel regime. In this project, each design decision is therefore treated as a testable hypothesis rather than as an assumed improvement.

Recent semantic communication studies also suggest that the motivation for this project should not be understood only as a replacement of conventional bit-level communication. A more accurate interpretation is that semantic communication complements conventional communication by introducing a new optimization layer above physical-layer reliability. In a conventional system, successful delivery is normally judged through bit error rate, symbol error rate, packet error rate, or throughput. These measures are still essential when the receiver must reproduce the exact transmitted message. However, in many intelligent communication scenarios, the receiver is not a passive data sink but a task-performing agent. For example, if a downstream classifier, dialogue agent, or control module can still make the correct decision from a compressed or partially corrupted representation, then exact bit-level reconstruction may be unnecessarily strict. This motivates a task-aware and meaning-aware design philosophy in which communication resources are allocated according to semantic importance rather than uniform syntactic fidelity.

This shift is especially relevant to future wireless systems because the expected traffic pattern is becoming increasingly machine-centric. Edge devices, autonomous agents, sensors, and embedded AI systems may exchange not only raw data but also intermediate features, intents, instructions, and task-specific semantic states. Under bandwidth limitation, latency constraint, and energy constraint, transmitting compact semantic representations can be more efficient than transmitting all raw symbols and reconstructing them perfectly. Therefore, a semantic transceiver should be evaluated not only by whether it can reproduce the sentence exactly, but also by whether it preserves enough meaning for the receiver to perform the intended reconstruction or task. This is consistent with the broader post-Shannon research trend, where semantic extraction, semantic encoding, adaptability, and reliability are treated as central challenges rather than secondary details \cite{liu2024survey_dcn, zhang2024semantic_nso}.

For this project, the most important motivation is robustness rather than compression alone. Text transmission through a neural semantic transceiver is vulnerable to channel noise in a different way from traditional coded communication. In a traditional coded system, errors are often localized at the bit or symbol level and can be analyzed through classical coding theory. In a neural semantic transceiver, however, channel perturbation affects continuous latent vectors. A small perturbation in the latent space may produce either a minor lexical variation or a large semantic shift after autoregressive decoding. This nonlinear error propagation makes low-SNR behavior especially important. The experimental focus on \SI{-10}{dB} and \SI{-5}{dB} therefore reflects the practical operating region in which robust semantic modelling is most strongly tested, not simply an arbitrary extension of the SNR grid.

Another motivation is the need to evaluate apparently reasonable neural-network improvements with caution. Semantic communication papers often report performance gains by adding attention modules, semantic losses, or pretrained language-model components. However, these modules may interact in unexpected ways. A semantic loss may encourage global sentence similarity but reduce token-level reconstruction. A language model may improve fluency but alter the original meaning. A channel-aware encoder may help representation learning but still provide insufficient information for decoder-side recovery. Therefore, each added module in this project is tested against a controlled baseline and judged by its measured effect under the same evaluation protocol.

\subsection{Research Questions and Objectives}

To make the investigation testable and reproducible, the project is organized around four explicit research questions that drive the design and experimental program:

\begin{enumerate}
    \item Can channel-quality-aware attention mechanisms, placed inside both the encoder and the decoder, improve semantic reconstruction quality under low SNR conditions?
    \item Does introducing a semantic-regularization term into the loss function consistently improve final communication quality, or does it instead introduce gradient conflict with token-level cross-entropy?
    \item Can a pretrained language model be used for semantic post-refinement in a way that provides reliable net gains on top of an already well-trained decoder?
    \item If learnt post-refinement proves unstable, can a training-free decoding strategy, such as beam search, offer a better engineering trade-off in terms of accuracy, cost, and determinism?
\end{enumerate}

Guided by these four questions, the concrete engineering objectives of the project are:

\begin{itemize}
    \item Build and verify a faithful baseline semantic communication transceiver for text transmission, compatible with the literature and reproducible on commodity hardware.
    \item Implement three distinct innovation modules covering model design, training objective, and inference refinement, each controlled by a clean command-line flag.
    \item Evaluate all variants under both AWGN and Rayleigh channels on an identical SNR grid and with identical metric definitions, so that every comparison is fair and well-scoped.
    \item Quantify effectiveness boundaries and failure mechanisms, not only best-case metric values, and report both positive and negative outcomes with clear attribution.
    \item Translate the observed evidence into concrete, deployment-oriented design recommendations that can guide follow-up final-year projects and early-stage research work.
\end{itemize}

To ensure that conclusions are objective, all model variants are evaluated under exactly the same data split, SNR grid, decoding termination criteria, and metric implementations. The project therefore explicitly emphasizes methodological fairness: when one module is added or removed, all other conditions remain fixed, and the resulting evaluation files are committed without manual cherry-picking. This principle is particularly important in semantic communication studies because several commonly used metrics, including BLEU, BERTScore, and sentence similarity, have different sensitivities to local and global errors, and inconsistent protocols could easily produce misleading conclusions in favor of a preferred variant.

\subsection{Contribution Summary}

The main contributions of this project can be summarized as follows:

\begin{itemize}
    \item A reproducible full-stack semantic communication pipeline is completed, including data preprocessing, model training, checkpointing, evaluation, visualization, and comparison scripts, with every step documented as command-line instructions in the appendix.
    \item The SNR-adaptive attention idea is extended from encoder-only usage, which is the dominant form reported in the literature, to decoder cross-attention, enabling explicit channel-conditioned memory retrieval at the reconstruction stage.
    \item A dynamic semantic-loss weighting mechanism is introduced to replace fixed weighting, so that the semantic-consistency term is de-emphasized under poor channel conditions and gradually re-emphasized under cleaner channels.
    \item A confidence-driven BERT masked-refinement strategy is implemented to replace earlier unconstrained GPT-style corrections, focusing edits on low-confidence positions and restricting global rewriting behavior.
    \item A beam search decoding strategy with beam width five is integrated and evaluated as an inference-stage method that requires no retraining.
    \item Negative experimental results are reported and analysed, including the weak performance of the semantic-loss component and the instability of LLM-based refinement.
\end{itemize}

In addition to the implemented modules, the project emphasizes controlled comparison across different design choices. Each major modification is evaluated under the same data split, SNR grid, channel setting, and metric definitions where applicable. This allows the effects of SNR-aware attention, semantic-loss weighting, LLM refinement, and beam search to be compared more directly. The resulting discussion therefore includes both performance improvements and negative results, rather than reporting only the best-performing configuration.

\subsection{Structure of This Report}

The remainder of this report is structured as follows. Section 2 presents the technical body, including literature context, the baseline architecture, the initial innovation hypotheses, the diagnosed limitations of those hypotheses, the resulting iterative redesign process, and the quantitative results of each configuration on AWGN and Rayleigh channels. Section 3 focuses on analysis and discussion, where the mechanisms behind the observed gains and failures are interpreted and connected with standard wireless-channel behavior and deep-learning optimization theory. Section 4 summarizes the conclusions and proposes concrete future work directions along both scientific and engineering axes. The references section then lists all external works cited in the main text, ordered by citation order.

The appendices provide reproducibility resources, including the complete command pipelines used for training and evaluation, and the mapping between experimental variants and saved evaluation files under the \texttt{checkpoints/} tree. Together, the main text and the appendices are designed to let a future student continue this line of work without reconstructing the entire experimental setup from scratch, and to allow examiners to verify any numerical claim by running the listed commands directly.

\FloatBarrier

\clearpage
\section{Body of the Report}
\subsection{Literature and Technical Background}

\subsubsection{From Shannon Communication to Semantic Communication}

Shannon's theory of communication provides the foundation of modern communication system design and formalizes the relationship between source entropy, channel capacity, and reliable transmission limits \cite{shannon1948, shannon1948part2}. However, Shannon's framework deliberately abstracts away the semantic meaning of the transmitted message in order to obtain clean analytical results, and treats information as a purely syntactic quantity. In contrast, the concept of semantic communication, traced back to the work of Carnap and Bar-Hillel \cite{carnap1954} and later formalized by Bao \emph{et al.} \cite{bao2011semantic}, explicitly emphasizes whether the received content preserves task-relevant meaning rather than whether each transmitted bit is perfectly recovered. Recent surveys \cite{lu2024survey, getu2024tutorial} describe semantic communication as a candidate enabler of post-Shannon communication systems, with applications spanning natural-language messaging, multimodal sensing, machine-to-machine coordination, and real-time control.

Deep learning has made semantic communication practically feasible by providing data-driven encoders that can compress text, speech, or image signals into latent representations that are robust to channel noise. A representative and highly influential milestone is DeepSC \cite{xie2021deepsc}, which demonstrated end-to-end semantic text transmission with Transformer-based encoders and decoders and introduced a joint source-channel coding view over neural latent. Subsequent work has extended semantic communication to speech recognition and synthesis \cite{shi2023speech}, to multimodal task-oriented transmission, and to wireless systems where source and channel coding are co-optimized with downstream task objectives.

The architectural foundation of these systems is closely related to the Transformer paradigm proposed by Vaswani \emph{et al.} \cite{vaswani2017attention}. Self-attention enables each token representation to aggregate global context in a single operation, which is especially valuable for preserving sentence-level semantic meaning when channel constraints force low-dimensional bottlenecks. Compared with recurrent models, Transformer-based transceivers provide better parallelism and cleaner modular decomposition, which is practical for iterative engineering experiments in which a single architectural component must be selectively replaced or ablated.

A useful way to position this project in the literature is to separate three related but distinct concepts: semantic-oriented communication, goal-oriented communication, and semantic-aware communication. Semantic-oriented communication focuses on transmitting meaning rather than exact symbols. Goal-oriented communication goes one step further and evaluates transmission according to whether the receiver can complete a downstream task. Semantic-aware communication keeps a conventional communication objective but uses semantic information to assist resource allocation, coding, or recovery. The present project mainly belongs to semantic-oriented communication because the primary output is still a reconstructed text sequence, but it also borrows ideas from goal-oriented communication by evaluating meaning-level preservation through BERTScore and sentence similarity rather than relying only on exact token match.

DeepSC is the most direct technical foundation for this project. It proposed a Transformer-based semantic communication system for text transmission, in which a semantic encoder extracts sentence-level meaning and a channel encoder maps the learned representation into channel symbols. The receiver then uses a channel decoder and semantic decoder to recover the transmitted sentence. This framework is attractive because it combines source coding and channel coding into a single trainable neural architecture, and because the Transformer is naturally suited to capturing long-range contextual dependencies in text \cite{xie2021deepsc}. However, this architecture also creates new design questions. If the semantic latent representation is corrupted by channel noise, the decoder must infer the most likely original text from a noisy continuous representation rather than from a clean symbolic sequence. This makes the receiver-side reconstruction strategy just as important as the transmitter-side semantic extraction.

The Transformer architecture itself is a key reason why semantic text communication became practical. Self-attention allows each token to attend to all other tokens in the sequence, which helps the encoder form contextual representations. Cross-attention then allows the decoder to condition each generated token on the received latent representation. Compared with recurrent sequence models, the Transformer is more parallelizable during training and more flexible for inserting conditioning signals such as SNR, channel-state information, or semantic importance weights \cite{vaswani2017attention}. This explains why the present project chooses Transformer modules as the baseline structure rather than a recurrent or convolutional text model.

Nevertheless, the literature also shows that semantic communication performance depends strongly on the channel model. Studies that evaluate both AWGN and Rayleigh fading commonly observe that Rayleigh fading is more challenging because it introduces multiplicative signal attenuation in addition to additive noise. This difference is important because an algorithm that performs well under AWGN may not automatically generalize to fading channels \cite{zhou2021semantic_adaptive_ut}. The present project's mixed Rayleigh results therefore should not be viewed as an implementation accident only; they are consistent with the broader observation that fading channels require richer adaptation than scalar-SNR conditioning.

\subsubsection{Why Channel Awareness Matters}

In wireless systems, channel quality directly affects the reliability of the features that arrive at the receiver. For deep transceivers, this observation has an important architectural implication: feature extraction and feature usage should both be channel-aware, not just one of them. Existing systems often condition the encoder on the current SNR, for example by concatenating an estimated SNR value with intermediate embeddings or by modulating attention weights with an SNR-derived gate. However, pure encoder-side conditioning is insufficient when the decoder must reconstruct under strong perturbation, because the decoder is the component that ultimately performs token-by-token recovery \cite{goldsmith2005wireless, rappaport2002wireless}.

Channel awareness also matters at the optimization level. If the training objective ignores channel statistics, the model may converge to a configuration that is optimal only for a narrow subset of the SNR distribution. This phenomenon is particularly severe under Rayleigh fading, where the effective SNR fluctuates sharply and deep fades occur with non-negligible probability. A well-designed channel-aware transceiver should therefore receive channel-quality signals at both encoding and decoding time and use them to dynamically adjust its internal behavior.

\subsubsection{Loss Design and Post-Refinement}

Cross-entropy remains the standard objective for sequence reconstruction in neural text communication and machine translation. To improve semantic fidelity, many studies add embedding-level or sentence-level regularization terms, so that two predictions with the same token accuracy can be further differentiated by how well they preserve sentence-level meaning. However, multi-objective optimization over token-level and semantic-level losses can cause gradient conflict, making nominally meaningful regularizers unstable in practice, and sometimes slowing convergence even when the long-run optimum improves \cite{kendall2018mtl, chen2018gradnorm}.

Post-decoding correction with pretrained language models is another common strategy that attempts to rescue decoded sequences by replacing low-probability or implausible tokens with more likely alternatives. However, objective mismatch between a generic language model and a communication reconstruction task may produce over-correction when the original decoded output is already strong \cite{devlin2019bert, sanh2019distilbert}. In other words, a language model that is trained to maximize linguistic plausibility is not necessarily aligned with the goal of preserving the transmitted content, and a single incorrect replacement can wipe out the marginal gain of several correct ones. This tension is a central motivation for the later confidence-driven design in this project.

In this project, output quality is measured through several complementary metrics rather than a single score. Word Accuracy quantifies how many tokens are recovered correctly and is directly sensitive to token-level noise. BLEU \cite{papineni2002bleu} provides n-gram overlap sensitivity across several n-gram orders and is widely used in machine translation. BERTScore \cite{zhang2020bertscore} captures contextual semantic similarity through pretrained representations. Sentence similarity, computed on sentence-level embeddings \cite{reimers2019sbert, gao2021simcse}, offers a global view of semantic alignment. Combining these metrics helps avoid one-dimensional interpretations, for example, a model that is fluent but semantically drifting, or a model that is semantically close but lexically unstable under low SNR.

The choice of evaluation metric is also a non-trivial design issue in semantic communication. In conventional digital communication, bit error rate and symbol error rate are natural because the receiver is expected to recover an exact bit stream. In semantic communication, these metrics are insufficient because two sentences may have different surface forms but nearly identical meanings. For example, ``the model improves at low SNR'' and ``low-SNR performance is improved by the model'' differ lexically but express similar content. Conversely, two sentences may share many words while differing in the key semantic relation. This creates a need for multiple complementary metrics.

BLEU is useful because it measures n-gram overlap and provides a stable, widely used lexical comparison. It is sensitive to local word order and phrase-level reconstruction, which makes it appropriate for evaluating whether the decoder preserves the original sentence structure \cite{papineni2002bleu}. However, BLEU alone can penalize valid paraphrases and may not fully capture contextual meaning. BERTScore addresses this limitation by comparing contextual embeddings from pretrained language models, so that semantically similar words or phrases can receive credit even when they are not exact string matches \cite{zhang2020bertscore}. In the context of this project, BLEU and Word Accuracy are therefore treated as reconstruction-oriented metrics, while BERTScore and sentence similarity are treated as meaning-oriented metrics.

This multi-metric design is important for interpreting the semantic-loss ablation. If the semantic-loss variant improved BERTScore while reducing BLEU-4 only slightly, the result might suggest a reasonable trade-off between lexical fidelity and semantic fidelity. However, the observed ablation shows that removing the semantic loss improves BLEU-4 at every SNR point, and the report does not show a strong compensating gain from the semantic-loss version. This means that the current semantic loss is not merely changing the metric preference; it is likely introducing an optimization conflict that weakens reconstruction quality. Therefore, the correct conclusion is not that semantic loss is useless, but that its present formulation is not sufficiently aligned with the actual decoding objective.

Post-refinement with a pretrained language model should be interpreted with similar caution. A language model is trained to produce fluent and probable text, while a communication receiver must recover faithful text. These objectives overlap but are not identical. When the decoded sentence is severely corrupted, a language model may help by replacing implausible words. When the decoded sentence is already mostly correct, the same model may introduce unnecessary substitutions. This explains why the project moved from unconstrained GPT-style rewriting to confidence-driven BERT masked refinement, and why even the constrained version remains risky. The problem is not only whether the language model is strong, but whether its editing objective is properly constrained by the communication task.

\subsection{Baseline System Design}

The baseline transceiver follows an Encoder--Channel--Decoder pipeline that is structurally aligned with DeepSC \cite{xie2021deepsc}. Input text is tokenized with a shared vocabulary, mapped to word embeddings with sinusoidal positional encoding, and then processed by a multi-layer Transformer encoder that generates semantic latent representations. These representations are projected into a channel space by a small channel encoder, passed through a differentiable channel model (either AWGN or Rayleigh), then reconstructed in sequence by a channel decoder and a Transformer decoder. The decoder is trained with teacher forcing and inference is performed in an autoregressive manner.

The baseline was selected because it reproduces the dominant semantic communication design pattern in recent literature while still allowing targeted module-level modifications. The baseline is trained with token-level cross-entropy and evaluated at fixed SNR points from \SI{-10}{dB} to \SI{20}{dB} in \SI{5}{dB} increments, so that downstream improvements can be attributed to specific added components rather than to data or training protocol changes. Training hyperparameters, such as optimizer type, learning rate schedule, sequence length cap, and batch size, are chosen to fit on a single-GPU workstation and are shared across all variants.

To make the implementation reproducible, the baseline and all improved variants share the same vocabulary, sequence length constraint, optimizer family, mixed-precision policy, and train/test split. This consistency prevents accidental gains from data leakage or from drifting in auxiliary components. It also makes ablation interpretation cleaner: if one metric changes after a module is swapped, the change can be traced to the intended module rather than to hidden side effects in preprocessing, tokenization, or evaluation. The preprocessing scripts and evaluation scripts are written to be deterministic when a random seed is provided, which further stabilizes comparisons.

The baseline is evaluated under both AWGN and Rayleigh fading channels because these two channels cover two practically important extremes: AWGN captures pure additive noise behavior without time variation, while Rayleigh introduces multiplicative fading, deep fades, and stronger statistical variability. Although real-world channels often combine additional effects, such as frequency selectivity, Doppler shifts, and interference, the AWGN--Rayleigh pair already provides a strong stress test for any semantic communication model and is the standard pairing in DeepSC-style studies.

\begin{figure}[!htb]
\centering
\ThesisIncludeGraphics{fig_architecture.png}
\caption{Block-level view of the semantic text transceiver (encoder, channel, decoder) with SNR-aware conditioning.}
\label{fig:architecture}
\end{figure}
Figure~\ref{fig:architecture} summarizes the end-to-end encoder–channel–decoder layout used throughout this report. It is placed here to connect the baseline description with the later modifications: SNR-aware attention is added inside the semantic encoder and decoder, while the channel encoder, wireless channel, and channel decoder define the communication bottleneck evaluated under AWGN and Rayleigh fading.

\subsection{Initial Innovation Hypotheses}

\subsubsection{Innovation A: SNR-Adaptive Attention}

The first hypothesis is that explicit SNR conditioning inside attention can improve robustness. If the model has access to an estimate of the current channel quality, it may adaptively balance semantic compactness and reconstruction detail. Under a very noisy channel, for example, the model should arguably favor more conservative attention patterns that aggregate context across a wider window and tolerate local corruption, whereas under a clean channel, sharper local attention can exploit fine-grained token dependencies. This first hypothesis predicts that an SNR-conditioned attention mechanism should outperform a non-adaptive one at the low-SNR end of the operating curve.

\subsubsection{Innovation B: Composite Semantic Loss}

The second hypothesis is that adding a semantic-consistency term to cross-entropy can improve meaning preservation:
\begin{equation}
L_{\mathrm{total}} = \alpha\, L_{\mathrm{ce}} + \beta\, L_{\mathrm{sem}},
\end{equation}
where $L_{\mathrm{ce}}$ is the token-level cross-entropy loss, $L_{\mathrm{sem}}$ is an embedding-based semantic consistency loss between encoder output and reconstructed decoder output, and $\alpha$, $\beta$ are non-negative weighting coefficients. In the initial formulation, $\beta$ was held fixed throughout training and across all SNR conditions, which made the composite loss easy to implement and analyze. This second hypothesis predicts that, for appropriately chosen $\beta$, the composite loss should outperform a pure cross-entropy loss on at least the high-SNR end of the operating curve, where the cross-entropy objective is close to saturation and additional semantic regularization should deliver marginal gains.

\subsubsection{Innovation C: LLM Post-Refinement}

The third hypothesis is that a pretrained language model can correct residual decoding errors through post-processing. In the first prototype, a compact DistilGPT-2 \cite{sanh2019distilbert} model was invoked to re-score and replace decoder output at the word level, driven by language-model probability only. The initial expectation was that language priors would restore semantically plausible words after noisy transmission, in particular by fixing spelling-like or grammar-like errors introduced by the channel.

At the hypothesis stage this idea appears intuitive, because language models encode strong grammatical and lexical priors learnt on very large corpora. However, communication reconstruction is not identical to generic text generation. The received message should remain faithful to the transmitted content, not merely plausible in natural language. This distinction is central to the failure analysis developed later in the project and ultimately motivates the transition from GPT-style free rewriting to BERT masked refinement driven by per-token confidence.

\subsection{Diagnosed Limitations and Root Causes}

After the first full experiment cycle, three key limitations were identified through a combination of metric inspection and manual sample inspection:
\begin{enumerate}
    \item Encoder-only SNR adaptation had limited transfer to more difficult fading scenarios. The encoder became channel-aware, but the decoder still had to reconstruct from perturbed features without any channel-quality signal.
    \item Fixed semantic-loss weighting produced weak negative transfer at several SNR points. The fixed weight was effectively too strong at low SNR, where the model needed to prioritize stable token recovery.
    \item External LLM post-refinement tended to over-correct and sometimes degraded final metrics, especially when the base decoding was already strong.
\end{enumerate}

A careful root-cause analysis suggests that channel uncertainty must be available at reconstruction time (not only at encoding time), that semantic supervision should be condition-dependent rather than flat, and that post-refinement should be constrained by per-token confidence and by task fidelity rather than by language-model probability alone. These observations motivated a transition from a “module addition” mindset to a “mechanism redesign” mindset. Instead of assuming that every semantically inspired component is beneficial, the project began treating each module as a hypothesis under falsifiable testing. This shift also made the subsequent redesign more evidence-based, because each module was retained, revised, or questioned according to its measured effect.

Manual inspection of incorrect samples provided additional insight. In many failure cases at \SI{-10}{dB}, the baseline decoder produced a short but almost correct sentence that was then rewritten by the LLM refiner into a longer, grammatically plausible, but semantically different sentence. Similarly, at high SNR, the fixed semantic loss occasionally biased the decoder toward paraphrased outputs that slightly reduced BLEU without obviously improving BERTScore. Both symptoms pointed toward the same underlying issue: unconditional objectives and unconstrained correction do not match the actual quality profile of a semantic communication system, which varies strongly with SNR.

\subsection{Iterative Redesign and Implementation}

\subsubsection{A': Decoder-side SNR-aware Cross-Attention}

SNR conditioning was extended from encoder self-attention to decoder cross-attention. The decoder now receives a channel-quality context vector while selecting memory features from the channel-decoded representations. This modification addresses the previously identified mismatch in which the encoder is channel-aware but the reconstructor still performs token generation without direct access to the current channel condition. In the modified cross-attention layer, the attention logits are additively biased by a lightweight neural module whose input contains the current SNR estimate and whose output provides head-specific attention modulation. This can be interpreted as introducing a learned prior over the decoder attention distribution conditioned on channel quality, allowing the model to dynamically adjust how strongly it aggregates received semantic features under different noise regimes. Implementation-wise, this change is localized to the decoder layer and does not alter the rest of the training, evaluation, or data-processing pipeline.


\begin{figure}[!htb]
\centering
\ThesisIncludeGraphics{fig_snr_attention.png}
\caption{SNR-adaptive attention: channel-quality information modulates self-attention in the encoder and cross-attention in the decoder (schematic).}
\label{fig:snr_attention}
\end{figure}
Figure~\ref{fig:snr_attention} illustrates SNR-conditioned attention in the encoder and decoder cross-attention paths. The key point is that SNR is not treated as an external evaluation variable only; it is injected into the attention computation so that feature aggregation can change with the estimated channel condition.

\subsubsection{B': Adaptive Semantic-Loss Weighting}

To reduce low-SNR instability caused by fixed semantic-loss weighting, the fixed $\beta$ was replaced by an SNR-dependent effective weight:
\begin{equation}
\beta_{\text{eff}} = \beta \cdot \sigma\!\left(\frac{\overline{\mathrm{SNR}}}{5.0}\right),
\end{equation}
where $\sigma$ is the sigmoid function and $\overline{\mathrm{SNR}}$ is the batch-average SNR in dB. This design weakens semantic regularization under poor channel conditions and gradually restores it under cleaner channels. The motivation is straightforward: when the channel is highly noisy, token-level reconstruction should dominate to stabilize optimization; when the channel is cleaner, the semantic regularizer can contribute more without being drowned in noise gradients. This conditional balancing is a lightweight alternative to more complex multi-objective solvers, such as GradNorm or uncertainty weighting \cite{kendall2018mtl, chen2018gradnorm}, and can be integrated without redesigning the overall training loop.

\begin{figure}[!htb]
\centering
\ThesisIncludeGraphics[0.80]{fig_composite_loss.png}
\caption{Composite training objective: cross-entropy $L_{\mathrm{ce}}$ and semantic term $L_{\mathrm{sem}}$ with effective weight $\beta_{\mathrm{eff}}$ controlled by batch SNR (conceptual diagram).}
\label{fig:composite_loss}
\end{figure}
Figure~\ref{fig:composite_loss} sketches the intended composite objective and the SNR-dependent scaling of the semantic branch. This diagram represents the design hypothesis tested in the ablation study, rather than a confirmed performance improvement.

\subsubsection{C': Confidence-driven BERT Masked Refinement}

The refinement stage was fundamentally redesigned so that the post-refiner no longer freely rewrites decoded sentences. The new workflow is:
\begin{enumerate}
    \item Obtain a per-position softmax confidence score directly from the decoder while decoding.
    \item Select a small number of positions whose confidence falls below an adaptive threshold.
    \item Apply masked replacement at those positions using a BERT masked-language-model head, restricted to the existing vocabulary and only accepting a replacement if the original and the candidate differ by at most one edit under a lexical-distance constraint.
\end{enumerate}

This design avoids global rewriting behavior and focuses corrections on positions where the decoder itself admitted uncertainty. In addition, hard constraints are introduced at the replacement stage, including a maximum number of allowed corrections per sentence and a minimum confidence gap before a replacement is accepted. Even with these constraints, experiments show that refinement remains sensitive to calibration quality of the confidence scores, which motivates further future work on confidence calibration.

\subsubsection{D: Beam Search Decoding}

A beam-search decoder with beam width five was added as an inference-stage enhancement. This change requires no retraining and serves as a controlled comparison against LLM-based refinement, since it improves output quality by exploring more hypotheses rather than by rewriting the top hypothesis. Beam search is included for two main reasons. First, it provides a deterministic and well-understood way to improve sequence-level hypothesis quality and is widely used in machine translation and speech recognition. Second, it acts as a strong baseline for “inference-only” enhancement, allowing a fair comparison with LLM refinement in terms of gain versus added complexity and latency.

Beam search can be interpreted as a receiver-side hypothesis management strategy. Greedy decoding selects the most probable token at each step, which is computationally cheap but locally optimal. Once a wrong token is selected early, later tokens are generated under an incorrect prefix, and the decoder may not be able to recover. This problem is especially severe at low SNR because the probability distribution over candidate tokens becomes flatter and more uncertain. Beam search reduces this risk by maintaining several partial hypotheses simultaneously and selecting the sequence with the best accumulated score after considering a wider search space.

In this project, beam search is particularly valuable because it improves the receiver without changing the trained model. This distinguishes it from architectural modifications and loss-function modifications, both of which require retraining and may introduce new optimization instability. Beam search therefore provides a clean test of whether the trained decoder already contains useful probability information that greedy decoding fails to exploit. The observed improvement at \SI{-10}{dB} supports this interpretation: the model has learned useful alternatives, but the greedy decoder sometimes commits too early to a suboptimal token.

Beam width five is chosen as a practical compromise. A very small beam may not differ much from greedy decoding, while a very large beam increases latency and can sometimes prefer overly generic or length-biased outputs. Since the project aims to provide an engineering recommendation rather than an exhaustive decoding study, beam width five is sufficient to test whether sequence-level search gives a meaningful benefit. Future work can extend this by sweeping beam widths from three to ten and measuring both performance and inference latency.

The comparison between beam search and BERT refinement is also methodologically important. Both are inference-stage enhancements, but they improve output quality through different mechanisms. Beam search searches within the probability distribution produced by the trained decoder, so it remains closely tied to the communication model. BERT refinement imports an external linguistic prior and edits the decoded sentence after the communication model has completed its work. The results show that the internal search strategy is more reliable than external editing in this setting. This supports the conclusion that receiver-side decoding should be optimized before adding a separate post-correction model.

Figure~\ref{fig:inference_pipeline} contrasts the three inference-stage alternatives evaluated in this project: greedy decoding, beam search, and confidence-driven BERT masked refinement. This figure clarifies why beam search and BERT refinement should be compared as separate inference-stage strategies rather than as equivalent correction mechanisms.
\begin{figure}[!htb]
\centering
\ThesisIncludeGraphics{fig_inference_pipeline.png}
\caption{Inference-stage alternatives: autoregressive decoding with confidence scores, beam search (width five), and optional BERT MLM refinement on low-confidence positions.}
\label{fig:inference_pipeline}
\end{figure}
\FloatBarrier

\subsection{Experimental Protocol and Evaluation Metrics}

All experiments use a unified protocol:

\begin{itemize}
    \item Channels: AWGN and Rayleigh.
    \item SNR points: $\{-10,-5,0,5,10,15,20\}$ dB.
    \item Metrics: Word Accuracy, BLEU-1/2/3/4, BERTScore, and Sentence Similarity.
    \item Decoding: greedy by default, beam search with beam width five where indicated.
    \item Refinement: none by default, confidence-driven BERT MLM where indicated.
    \item Optimizer: AdamW \cite{loshchilov2019adamw} with a warmup-then-decay learning-rate schedule, dropout fixed across variants.
\end{itemize}


\begin{figure}[!htb]
\centering
\ThesisIncludeGraphics{fig_channel_models.pdf}
\caption{Conceptual illustration of AWGN (scatter around a transmitted symbol) and slow Rayleigh-like fading (time-varying envelope). Qualitative only.}
\label{fig:channel_models}
\end{figure}
Figure~\ref{fig:channel_models} gives a conceptual view of the two channel models used in simulation. It is included to clarify why AWGN and Rayleigh should not be treated as equivalent noise conditions: AWGN mainly introduces additive perturbation, whereas Rayleigh fading also changes the received signal amplitude over time.


\begin{figure}[!htb]
\centering
\ThesisIncludeGraphics{fig_workflow.png}
\caption{Experimental workflow: data preparation, training of baseline and improved variants, ablation, evaluation over the SNR grid, and saving of JSON metrics.}
\label{fig:workflow}
\end{figure}
Figure~\ref{fig:workflow} outlines the experimental workflow from preprocessing to evaluation exports.

This workflow also explains how fair comparison is maintained: each variant is trained or evaluated within the same preprocessing, SNR grid, and metric-export pipeline.

The evaluated model groups are:
\begin{enumerate}
    \item Baseline-AWGN.
    \item Improved-AWGN.
    \item Improved-AWGN-NoSem (ablation of semantic loss).
    \item Baseline-Rayleigh.
    \item Improved-Rayleigh.
    \item Improved-AWGN + Beam5.
    \item Improved-AWGN + FT-BERT (BERT fine-tuned on Europarl, used as refiner).
    \item Improved-AWGN + Beam5 + FT-BERT.
\end{enumerate}

All evaluations share the same sample cap, tokenization, and decoding termination settings to ensure fairness. For LLM-enabled runs, the confidence collection procedure and the word-level confidence mapping are fixed and deterministic, so that the same decoder output always triggers the same refinement decisions. This removes a class of stochastic variance that could otherwise obscure small but consistent differences between variants. The evaluation results are written to variant-specific JSON files under the \texttt{checkpoints/} tree, so that every reported number can be re-computed from the saved outputs rather than being regenerated at reporting time.

Beyond these quantitative metrics, a small number of qualitative samples are also collected at each SNR point, covering representative clean, noisy, and failed cases. These samples support manual inspection and are used to support mechanism-level discussion in Section 3. The combined use of quantitative metrics and qualitative samples provides a more complete picture of system behavior than any single aggregate number could deliver.

\FloatBarrier
\subsection{Result Presentation}

\subsubsection{AWGN Main Comparison}
\begin{table}[htbp]
\centering
\caption{Main AWGN comparison (Baseline vs Improved, greedy decoding).}
\begin{tabular}{cccccc}
\toprule
SNR (dB) & \multicolumn{2}{c}{WordAcc} & \multicolumn{2}{c}{BLEU-4} & BLEU-4 Gain \\
\cmidrule(lr){2-3}\cmidrule(lr){4-5}
& Baseline & Improved & Baseline & Improved & (Improved - Baseline) \\
\midrule
-10 & 0.4425 & 0.4773 & 0.4189 & 0.4329 & +0.0140 \\
-5  & 0.5841 & 0.6130 & 0.7340 & 0.7386 & +0.0046 \\
0   & 0.7325 & 0.7440 & 0.8371 & 0.8375 & +0.0004 \\
5   & 0.9567 & 0.9674 & 0.9633 & 0.9685 & +0.0052 \\
10  & 0.9941 & 0.9954 & 0.9912 & 0.9934 & +0.0022 \\
15  & 0.9956 & 0.9965 & 0.9931 & 0.9944 & +0.0013 \\
20  & 0.9961 & 0.9966 & 0.9934 & 0.9945 & +0.0011 \\
\bottomrule
\end{tabular}
\label{tab:awgn-main}
\end{table}



Table~\ref{tab:awgn-main} compares the Baseline and Improved transceivers under AWGN with greedy decoding. The table shows that improvements are concentrated in low-SNR ranges, where robust representation learning is most critical. At high SNR, absolute scores are already close to their practical ceiling, so additional module gains naturally shrink and the two curves become nearly indistinguishable. The most visible gain appears at \SI{-10}{dB}, where BLEU-4 increases from 0.4189 to 0.4329 and Word Accuracy increases from 0.4425 to 0.4773. These gains are consistent with the design intent of the channel-aware redesign.

\subsubsection{Ablation: Semantic Loss}
\begin{table}[htbp]
\centering
\caption{Ablation of semantic loss under AWGN (Improved vs Improved-NoSem).}
\begin{tabular}{cccccc}
\toprule
SNR (dB) & \multicolumn{2}{c}{WordAcc} & \multicolumn{2}{c}{BLEU-4} & NoSem - Improved (BLEU-4) \\
\cmidrule(lr){2-3}\cmidrule(lr){4-5}
& Improved & NoSem & Improved & NoSem & \\
\midrule
-10 & 0.4773 & 0.4837 & 0.4329 & 0.4488 & +0.0159 \\
-5  & 0.6130 & 0.6210 & 0.7386 & 0.7448 & +0.0062 \\
0   & 0.7440 & 0.7446 & 0.8375 & 0.8398 & +0.0023 \\
5   & 0.9674 & 0.9683 & 0.9685 & 0.9709 & +0.0024 \\
10  & 0.9954 & 0.9964 & 0.9934 & 0.9946 & +0.0012 \\
15  & 0.9965 & 0.9976 & 0.9944 & 0.9959 & +0.0015 \\
20  & 0.9966 & 0.9976 & 0.9945 & 0.9960 & +0.0016 \\
\bottomrule
\end{tabular}
\label{tab:ablation}
\end{table}



Table~\ref{tab:ablation} compares the Improved configuration with a matching configuration in which the semantic-loss term is disabled during training (NoSem). The NoSem variant outperforms the semantic-loss variant at every tested SNR point. This indicates that the semantic regularizer, as currently formulated, is not yet aligned with the downstream reconstruction objective. This negative result indicates that a theoretically motivated component can still degrade performance when the multi-objective training signal is not well calibrated. The ablation therefore flags semantic loss as a module that requires further redesign rather than outright removal.

\subsubsection{Cross-Channel Comparison (Rayleigh)}
\begin{table}[htbp]
\centering
\caption{Rayleigh channel comparison (Baseline vs Improved).}
\begin{tabular}{cccccc}
\toprule
SNR (dB) & \multicolumn{2}{c}{WordAcc} & \multicolumn{2}{c}{BLEU-4} & BLEU-4 Gain \\
\cmidrule(lr){2-3}\cmidrule(lr){4-5}
& Baseline & Improved & Baseline & Improved & (Improved - Baseline) \\
\midrule
-10 & 0.4730 & 0.4739 & 0.4697 & 0.4640 & -0.0057 \\
-5  & 0.6044 & 0.6014 & 0.6662 & 0.6603 & -0.0059 \\
0   & 0.7501 & 0.7527 & 0.8170 & 0.8055 & -0.0115 \\
5   & 0.8900 & 0.8938 & 0.9167 & 0.9148 & -0.0019 \\
10  & 0.9550 & 0.9572 & 0.9646 & 0.9647 & +0.0001 \\
15  & 0.9824 & 0.9840 & 0.9825 & 0.9828 & +0.0003 \\
20  & 0.9886 & 0.9915 & 0.9869 & 0.9876 & +0.0007 \\
\bottomrule
\end{tabular}
\label{tab:rayleigh}
\end{table}



Table~\ref{tab:rayleigh} reports the Rayleigh counterpart of the main AWGN comparison. Improvements are limited and mixed across the SNR range, and in several low and medium SNR points the Improved variant is slightly worse than Baseline. Only at high SNR does the difference become near-zero or weakly positive. This suggests that channel-aware design alone is insufficient under strong fading variability and that future work must integrate richer channel-state cues than a single scalar SNR.

\subsubsection{Decoding and Refinement Strategy Comparison}
\begin{table}[htbp]
\centering
\caption{Inference strategy comparison on Improved-AWGN (BLEU-4).}
\begin{tabular}{ccccc}
\toprule
SNR (dB) & Greedy (No LLM) & Beam5 & FT-BERT (LLM) & Beam5 + FT-BERT \\
\midrule
-10 & 0.4329 & 0.4552 & 0.4191 & 0.4046 \\
-5  & 0.7386 & 0.7496 & 0.7348 & 0.7326 \\
0   & 0.8375 & 0.8494 & 0.8339 & 0.8341 \\
5   & 0.9685 & 0.9726 & 0.9650 & 0.9654 \\
10  & 0.9934 & 0.9934 & 0.9905 & 0.9910 \\
15  & 0.9944 & 0.9945 & 0.9934 & 0.9935 \\
20  & 0.9945 & 0.9946 & 0.9933 & 0.9938 \\
\bottomrule
\end{tabular}
\label{tab:strategy}
\end{table}



Finally, Table~\ref{tab:strategy} compares decoding and refinement strategies on Improved-AWGN. Beam-only consistently beats greedy decoding at low SNR, while adding FT-BERT on top of beam reduces performance. This interaction effect is crucial: stronger base decoding reduces the remaining room for post-editing and increases the probability that external correction introduces errors rather than removing them.

\FloatBarrier
\subsection{Interim Summary}

Taken together, the quantitative results reveal three main observations. First, the channel-aware redesign improves low-SNR AWGN performance, which is consistent with the original hypothesis. Second, semantic-loss regularization remains unstable in the current formulation, and the ablation results show that it reduces BLEU-4 across the tested SNR range. Third, beam search is a more reliable inference-stage enhancement than external LLM refinement in this experimental setting, particularly when the base decoder already performs well.

These observations indicate that the final system should not simply include every proposed module. Instead, the evidence supports retaining components with consistent measured gains, further revising components with mixed effects, and treating harmful components as limitations for future investigation. The following section analyses the mechanisms behind these results before the conclusions section summarizes the final design recommendations.

\FloatBarrier

\clearpage
\section{Analysis and Discussion}
\subsection{Main Findings under AWGN}

The AWGN results in Table~\ref{tab:awgn-main} show a clear and interpretable pattern: channel-aware improvements are most effective in challenging SNR regimes and gradually saturate at high SNR. At \SI{-10}{dB}, the Improved model yields a BLEU-4 gain of +0.0140 and a Word Accuracy gain of +0.0348 over the Baseline. Word Accuracy gain is particularly meaningful because it reflects direct token reconstruction quality under severe channel noise, where any systematic improvement necessarily propagates into subsequent n-gram metrics and into sentence-level semantic metrics. The gain is also consistent with the design motivation of the decoder-side SNR-aware cross-attention, which gives the reconstructor explicit access to channel quality at the moment of token generation.

Figures~\ref{fig:bleu_awgn} and~\ref{fig:wordacc_awgn} visualize the AWGN trends corresponding to Table~\ref{tab:awgn-main}: Figure~\ref{fig:bleu_awgn} focuses on BLEU-$4$, while Figure~\ref{fig:wordacc_awgn} shows Word Accuracy, making the low-SNR gain and high-SNR saturation easier to see.
\begin{figure}[!htb]
\centering
\ThesisIncludeGraphics{fig_bleu_awgn_strategies.pdf}
\caption{AWGN: BLEU-$4$ versus SNR for baseline and improved configurations (greedy, Beam5, and FT-BERT post-processing where \texttt{LLM\_BLEU-4} is present in the export).}
\label{fig:bleu_awgn}
\end{figure}
\begin{figure}[!htb]
\centering
\ThesisIncludeGraphics{fig_wordacc_awgn.pdf}
\caption{AWGN: word accuracy versus SNR for baseline versus improved (greedy decoding).}
\label{fig:wordacc_awgn}
\end{figure}
\FloatBarrier

At medium and high SNR (from \SI{5}{dB} upward), both models approach very high absolute scores (BLEU-4 above 0.96 and Word Accuracy above 0.95), so relative gains become small. This saturation does not invalidate the improvement. Instead, it confirms that low-SNR robustness is the true stress test for semantic communication transceivers, and that gains measured in a ``safe'' regime should not be interpreted as evidence of general superiority. Reporting both the low-SNR and high-SNR halves of the curve is therefore essential and is the reason why the experimental grid is extended down to \SI{-10}{dB} rather than stopping at \SI{0}{dB}.

Another practical observation is metric consistency at low SNR: BLEU-4, Word Accuracy, BERTScore, and Sentence Similarity all move in the same direction for the Improved model relative to the Baseline at \SI{-10}{dB}. This joint movement increases the confidence that the gains are not metric artefacts. In many text-generation tasks, lexical and semantic metrics may diverge sharply, so the fact that they agree here supports the claim that channel-aware redesign improves both local token accuracy and global sentence-level reconstruction quality. The joint movement also implies that the Improved model does not simply trade BLEU for paraphrase-like semantic matches, which would be a typical failure mode of semantic-only regularization.

Qualitative inspection of representative low-SNR outputs also suggests that early decoding errors are important. When the first few tokens are corrupted, later predictions are often affected by the incorrect prefix. This is consistent with the autoregressive nature of the decoder, where each prediction depends on previously generated tokens. Therefore, improvements in early-token reliability can have a noticeable effect on the overall sequence quality, especially under severe channel noise.

\FloatBarrier
\subsection{Ablation Analysis: Semantic Loss}

The ablation results in Table~\ref{tab:ablation} reveal an important negative finding: removing the semantic-loss term improves BLEU-4 at every tested SNR point. The largest gap appears at \SI{-10}{dB}, where the NoSem variant exceeds the semantic-loss version by +0.0159 in BLEU-4. This indicates that, under the current implementation, semantic regularization introduces weak but consistent negative transfer. The result is particularly important because semantic communication is often expected to benefit naturally from a semantic-level objective. However, the evidence shows that a conceptually reasonable auxiliary loss does not automatically improve the final communication objective.

The most direct explanation is gradient conflict between token-level reconstruction and embedding-level semantic alignment. Cross-entropy encourages the decoder to assign high probability to the exact target token at each autoregressive step. In contrast, the semantic-consistency loss encourages the reconstructed representation to remain close to a reference embedding. These objectives may agree in many cases, but they are not guaranteed to point in the same optimization direction. For example, if the target word is “increase” and the decoder predicts “improve”, the two words may be close in semantic embedding space, but the prediction is still incorrect from the perspective of token-level reconstruction. In such a case, the semantic loss may partially reward a prediction that cross-entropy must penalize, making the composite training signal less focused.

This conflict becomes more harmful under low-SNR conditions. When the channel perturbation is strong, the received latent representation is already noisy and the decoder must prioritize stable token recovery. If the semantic-loss weight remains non-zero in this regime, the model may be encouraged to optimize a noisy or weakly discriminative semantic target instead of learning a sharper reconstruction boundary. The sigmoid-based SNR-conditioned weighting partly addresses this issue by reducing the semantic-loss contribution at low SNR, but the ablation results show that this adjustment is not sufficient to remove the negative effect.

The fact that the NoSem variant outperforms the semantic-loss variant across all SNR points also suggests that the problem is not only low-SNR over-regularization. If the semantic loss were merely too strong under noisy conditions, one would expect the semantic-loss version to recover or outperform the NoSem version at high SNR, where the received representation is cleaner. Since this does not occur, the more likely explanation is objective misalignment: the specific semantic-consistency term used in this project does not provide useful additional supervision for the final autoregressive reconstruction task.

Therefore, Innovation B should be interpreted as a valuable but unsuccessful attempt rather than a confirmed improvement. The project does not show that semantic regularization is useless in general. Instead, it shows that the present formulation is not sufficiently aligned with the downstream reconstruction objective. Future work should redesign this component using conflict-aware weighting, delayed semantic-loss activation, uncertainty-based loss scaling, or contrastive sentence-level objectives. Such methods may preserve the intended meaning-level supervision while reducing early-stage interference with token-level learning.

\FloatBarrier
\subsection{Cross-Channel Analysis: Rayleigh Robustness}

While the AWGN results validate the effectiveness of the proposed design under additive noise, it is equally important to examine whether these gains generalize to more realistic fading channels. Table~\ref{tab:rayleigh} shows that AWGN-effective improvements do not directly transfer to Rayleigh fading. In the low and medium SNR ranges, BLEU-4 is slightly lower for the Improved model than for the Baseline, and only at high SNR does the difference become near-zero or weakly positive. This behavior is notable because the Improved model was explicitly designed to be channel-aware, yet it does not consistently outperform the Baseline under the more variable fading channel.

The main reason is that scalar SNR conditioning captures only average channel quality, not the instantaneous structure of a fading process. In an AWGN channel, the transmitted symbols are mainly corrupted by independent additive Gaussian noise. Under this assumption, scalar SNR is a relatively informative summary of the channel condition. In a Rayleigh fading channel, however, the received symbols are affected by multiplicative fading coefficients as well as additive noise. Even when the average SNR is fixed, the reliability of individual transmitted symbols can vary substantially due to deep fades and short-time channel fluctuation.

From a representation perspective, Rayleigh fading can distort the semantic latent space in a less uniform way than AWGN. AWGN tends to perturb the transmitted latent representation around its original location, whereas fading may suppress or amplify different parts of the transmitted signal before additive noise is introduced. If the decoder only receives an average SNR value, it may know that the channel is generally poor, but it cannot identify which latent dimensions or symbol positions are particularly unreliable. Therefore, the decoder-side SNR-aware attention can adjust its global behavior, but it cannot selectively distrust corrupted components of the received representation.

This explains why “SNR-aware” is not equivalent to “fading-aware”. The current design gives the model useful information about average channel severity, but it does not provide sufficiently expressive information about fading dynamics. A more robust Rayleigh-oriented design would require richer channel-state descriptors, such as estimated channel coefficients, short-term fading statistics, per-symbol reliability indicators, or learned channel-state embeddings. These features could allow the receiver to distinguish between globally noisy transmission and locally unreliable latent symbols.

Another plausible factor is the training distribution. If the Rayleigh training schedule does not sufficiently emphasize low-SNR and deep-fade cases, the model may allocate too much capacity to relatively easier conditions and fail to develop strong robustness against severe fading events. A curriculum-style training strategy that progressively increases exposure to difficult fading samples could therefore improve generalization. Overall, the key limitation is not the absence of channel awareness, but the insufficient expressiveness of scalar SNR as a descriptor of fading dynamics.

Figure~\ref{fig:bleu_rayleigh} plots BLEU-$4$ versus SNR for the Rayleigh-trained baseline and improved models (greedy decoding).
\begin{figure}[!htb]
\centering
\ThesisIncludeGraphics{fig_bleu_rayleigh.pdf}
\caption{Rayleigh channel: BLEU-$4$ versus SNR for baseline and improved transceivers (greedy decoding).}
\label{fig:bleu_rayleigh}
\end{figure}
\FloatBarrier

\subsection{LLM Refinement Analysis}

The FT-BERT results in Table~\ref{tab:strategy} remain below the no-LLM greedy baseline across all SNR points, even after the refiner has been fine-tuned on the Europarl corpus. Although this approach is much more controlled than the earlier unconstrained GPT-style rewriting, it still introduces harmful modifications when the base decoder is already performing reasonably well. This outcome is consistent with the broader narrative of this project: semantic communication reconstruction is not identical to natural language modelling, and improving linguistic plausibility is not the same as improving transmission fidelity.

The most likely mechanism is \textbf{over-correction under low error entropy}. When base decoding quality is already strong, even a small number of wrong replacements can dominate marginal improvements. In other words, language plausibility and transmission fidelity are related but not identical objectives, and a replacement that improves plausibility can still damage fidelity. The confidence-driven masking partially mitigates this issue but cannot eliminate it, because a low-confidence decoder token can still happen to be correct, and BERT MLM has no way to see the original transmitted message.

A second contributing factor is domain mismatch. Although the refiner is fine-tuned on Europarl and is therefore closer to the evaluation domain than an off-the-shelf BERT, fine-tuning mainly adjusts the vocabulary distribution rather than the conditional distribution over possible corrections of noisy decoder output. A more effective refinement would require confidence scores that are actually calibrated against true token errors, which neither the decoder nor the BERT head provides out of the box. This is the rationale for the Future Work proposal on dedicated confidence calibration.

A third observation is that the combined Beam5 + FT-BERT configuration is the worst of all tested configurations at \SI{-10}{dB}. This is not a coincidence: once beam search has already absorbed most of the recoverable decoding errors, additional post-editing operates in a “low-entropy” regime where virtually every correction is a coin flip. This again supports the ordering principle developed in the conclusions section, which is to maximize base decoder quality first and only apply external correction when the decoder is demonstrably uncertain, and the candidate replacement can be validated under a fidelity constraint.

\subsection{Decoding Strategy Analysis}

Among the post-training modifications evaluated in this project, beam search gives the most consistent improvement in the tested low-SNR AWGN setting. At \SI{-10}{dB}, BLEU-4 rises from 0.4329 under greedy decoding to 0.4552 under Beam5, giving an absolute gain of +0.0223. Gains also remain visible at \SI{-5}{dB} and \SI{0}{dB}, while they become negligible at higher SNR values because the decoder is already close to the practical performance ceiling.

The reason for this improvement can be understood from a probabilistic perspective. Greedy decoding selects the most probable token at each step and therefore optimizes locally. Once an incorrect token is selected early, all later predictions are conditioned on a wrong prefix, and the decoder may not be able to recover. This problem is especially severe at low SNR because channel noise flattens the token probability distribution and increases uncertainty among competing candidates. Beam search better approximates the sequence-level maximum-likelihood solution by maintaining several partial hypotheses and selecting the best complete sequence after considering a wider search space.

This also explains why beam search is particularly suitable for semantic communication under noisy channels. The received latent representation may still contain useful information even when the top token at a single decoding step is unreliable. Greedy decoding may discard that information too early, while beam search preserves alternative hypotheses long enough for later context to disambiguate them. Therefore, beam search does not create new semantic knowledge; it makes better use of the probability distribution already learned by the trained decoder.

The comparison with FT-BERT refinement further clarifies the advantage of beam search. Both methods are inference-stage enhancements, but they operate through different mechanisms. Beam search explores alternative sequences within the decoder’s own probability model, so its decisions remain tied to the received latent representation. In contrast, BERT refinement introduces an external masked-language model that mainly evaluates linguistic plausibility after decoding has already occurred. Although this external prior can be useful when the decoded sentence is severely corrupted, it can also replace a correct but low-confidence token with a fluent but incorrect alternative.

The combined Beam5 + FT-BERT configuration performs worse than Beam5 alone at every tested SNR point. This indicates that once beam search has already improved the base hypothesis, the remaining errors become too sparse, too content-specific, or too uncertain for generic masked-language correction to fix reliably. In this low-error regime, external refinement has limited positive headroom but substantial downside risk. A single unnecessary replacement can reduce BLEU, Word Accuracy, and semantic fidelity, especially if the replacement changes a key content word.

From an engineering perspective, beam search provides a useful trade-off in this experimental setting because it requires no retraining, introduces no additional language-model dependency, and improves the low-SNR AWGN cases tested in this project. The measured latency increase of approximately 1.8 times compared with greedy decoding is acceptable for the text-based scenario considered here, although this cost should still be reassessed for longer sequences or real-time applications. Therefore, Beam5 is a reasonable inference choice for the current system, while LLM-based refinement should remain conditional on decoder uncertainty and should require additional fidelity constraints before deployment.

\subsection{Integrated Design Insights}

By combining all available evidence, the project proposes three practical design principles for future semantic communication transceivers:

\begin{enumerate}
    \item \textbf{Prioritize in-model channel awareness} over external correction. Channel-quality information should reach both encoding and decoding stages, because downstream reconstruction carries most of the reliability burden.
    \item \textbf{Treat semantic regularization as conditional}, not universally beneficial. Fixed or naively scheduled semantic losses can produce negative transfer. Future implementations should rely on conflict-aware weighting or gated activation rather than on fixed hyperparameters.
    \item \textbf{Prefer conservative inference upgrades} (beam search) before introducing flexible post-editing. Inference strategy gains are reproducible, low-risk, and require no retraining, whereas LLM post-refinement can act as a silent regression when base quality is high.
\end{enumerate}

These observations are based on the measured results rather than on the initial design assumptions. They show that the usefulness of a module depends on its interaction with the channel model, training objective, and decoding procedure. In particular, modules that appear reasonable in isolation may become ineffective or harmful after being combined with other components.

The main implication is that semantic communication transceiver design should be evaluated through controlled comparisons rather than through architectural complexity alone. The results connect metric changes to specific mechanisms, including decoder-side channel awareness, gradient conflict in multi-objective training, confidence calibration in post-refinement, and sequence-level uncertainty during decoding.

\subsection{Discussion Summary}

The discussion identifies different levels of support for the proposed modifications. Innovation A is supported under low-SNR AWGN conditions, where the Improved model outperforms the Baseline in BLEU-4 and Word Accuracy. Innovation B is not supported by the current ablation results, since removing the semantic-loss term improves BLEU-4 across all tested SNR points. Innovation C also remains uncertain, because confidence-driven FT-BERT refinement does not provide reliable net gains on top of the Improved decoder.

Beam search gives the clearest inference-stage improvement in the current experiments. Unlike external refinement, it does not introduce an additional language model or rewrite the decoded sentence after the communication process. Instead, it improves sequence selection within the decoder’s own probability distribution. This makes it less likely to introduce semantic drift than LLM-based post-refinement.

Overall, the results suggest that robustness improvements in semantic communication are strongly dependent on channel regime, decoder quality, objective alignment, and inference-time interaction. Therefore, each module should be evaluated not only by its standalone motivation, but also by its measured effect after integration into the full transceiver.

\FloatBarrier

\clearpage
\section{Conclusions and Further Work}
\subsection{Conclusions}

This project investigated the improvement of a DeepSC-style semantic transceiver through baseline comparison, module-level ablation, cross-channel evaluation, and inference-strategy analysis. The same evaluation protocol was used across the main model variants so that the effects of different components could be compared under consistent conditions. The conclusions below are based on the quantitative results reported in Section 2 and the mechanism-level analysis in Section 3.

The key conclusions that emerge from the quantitative evidence are:
\begin{enumerate}
    \item \textbf{SNR-adaptive transceiver redesign is effective under AWGN low-SNR conditions.} The Improved model clearly outperforms the Baseline at \SI{-10}{dB} and \SI{-5}{dB}, with joint gains across BLEU-4 and Word Accuracy, which reduces the likelihood that the improvement is a metric artefact.
    \item \textbf{Semantic-loss regularization is not yet stable in this implementation.} The ablation study shows that removing the semantic term improves final BLEU-4 at every tested SNR point, indicating that the current fixed and SNR-scaled formulations are not well aligned with the reconstruction objective.
    \item \textbf{BERT-based confidence refinement does not produce robust net gains.} The FT-BERT refiner mitigates some severe failures compared with earlier unconstrained GPT-style corrections but still tends to over-correct when the base decoder is strong, and it is consistently beaten by the no-refinement configuration.
    \item \textbf{Beam search is the best practical enhancement in this project.} It yields consistent low-SNR gains without retraining, and its gains dominate those of LLM-based refinement in every head-to-head comparison.
    \item \textbf{Rayleigh robustness remains an open challenge.} Improvements that are effective in AWGN do not transfer directly to fading-dominant conditions, which points to a genuine gap between scalar-SNR conditioning and richer fading-state conditioning.
\end{enumerate}

From an engineering perspective, the recommended deployment configuration for the text-based semantic communication scenario studied in this project is therefore: \textbf{Improved transceiver + Beam5 decoding, without external LLM refinement}. This recommendation is grounded in the balance between effectiveness and implementation cost. The Improved transceiver introduces manageable architectural changes relative to the baseline and provides meaningful low-SNR gains. Beam search adds a modest decoding-time cost but no retraining overhead and consistently improves the most difficult cases. In contrast, external LLM refinement increases dependency complexity while offering unstable net benefit in the current setup and is therefore not recommended in its current form.

\subsection{Suggestions for Further Work}

Future work should focus not only on improving absolute performance, but also on resolving the specific limitations identified by the experiments. The results suggest that the most important next steps are richer fading-aware channel modelling, conflict-aware semantic objective design, communication-constrained post-refinement, and complexity-aware deployment evaluation. These directions are directly motivated by the observed failure modes rather than added as generic extensions.

\begin{enumerate}
    \item \textbf{Rayleigh-specific adaptation:} incorporate richer channel-state features beyond scalar SNR, such as fading statistics, short-time variance, or estimated CSI embeddings. A natural next step is to train the model with curriculum-style sampling that emphasizes deep-fade events.
    \item \textbf{Conflict-aware multi-objective training:} apply gradient balancing methods such as GradNorm-style strategies or uncertainty weighting to reduce CE--semantic objective interference \cite{kendall2018mtl, chen2018gradnorm}. Alternatively, investigate delayed activation of the semantic loss after the cross-entropy loss has converged to a reasonable level.
    \item \textbf{Constrained post-refinement:} replace free replacement with edit-distance constraints, lexical whitelist constraints, or N-best reranking over the beam output. A hybrid scheme that only accepts a replacement when it improves a calibrated confidence signal and does not decrease beam score could address the current over-correction failure.
    \item \textbf{Confidence calibration:} jointly learn decoder confidence quality so that low-confidence triggers align better with true token errors. Temperature scaling, focal-loss fine-tuning, or ensemble-based calibration are concrete technical candidates.
    \item \textbf{Domain and task expansion:} validate on additional corpora (for example WMT or news commentary) and on task-oriented semantic communication settings (question answering, intent classification) to test the generalization of the findings.
    \item \textbf{End-to-end uncertainty modelling:} jointly model channel and decoding uncertainty so that both decoding and refinement decisions can be driven by a single consistent uncertainty signal.
    \item \textbf{Beam width sweep and deployment profiling:} investigate beam widths between 3 and 10, quantify the corresponding latency and memory cost, and produce Pareto curves that inform deployment trade-offs.
\end{enumerate}

Communication-aware LLM refinement: future work should redesign LLM refinement as a communication-constrained reranking problem rather than a free editing problem. Instead of allowing BERT or another language model to directly replace tokens, the system could generate a small candidate set from beam search and then use the language model only as one component of a combined score. The final score could include decoder log probability, semantic similarity to the received latent representation, language-model fluency, and a penalty for unnecessary edits. This would preserve the benefit of pretrained linguistic knowledge while reducing the risk of semantic drift.
Channel-state-aware attention: the current SNR-adaptive attention mechanism uses scalar SNR as the main channel-quality indicator. This is sufficient for a first design but insufficient for fading-dominant channels. A future system could feed the decoder with estimated channel-state information, per-symbol reliability, or a learned channel embedding. The attention module could then assign lower weights to unreliable latent positions and higher weights to stable positions. This would make the decoder more selective under Rayleigh fading and may directly address the weak cross-channel generalization observed in this project.
Statistical significance and repeated-seed validation: the current report uses controlled single-run comparisons with identical data splits and evaluation settings. This is appropriate for an undergraduate engineering project, but future research should repeat each experiment with multiple random seeds. Mean values, standard deviations, and confidence intervals should be reported for BLEU-4, Word Accuracy, BERTScore, and sentence similarity. This would make small differences at high SNR easier to interpret and would prevent over-interpreting marginal gains.
Task-oriented evaluation: future work should test whether the reconstructed sentence is sufficient for a downstream task, such as intent classification, question answering, or command execution. This would move the system closer to goal-oriented semantic communication. In such a setting, exact sentence reconstruction may be less important than whether the receiver makes the correct task decision. A task-level metric could therefore reveal benefits that BLEU or Word Accuracy may miss.
Complexity-aware deployment study: the final recommended system uses the improved transceiver with Beam5 decoding. Although Beam5 requires no retraining, it increases inference cost relative to greedy decoding. Future work should measure latency, GPU memory usage, CPU inference time, and energy consumption. The final output should be a Pareto comparison between reconstruction quality and computational cost. This would make the project more relevant to edge-device deployment, where semantic communication is often expected to operate under limited hardware resources.
Robustness to domain shift: the current text data are sufficient for controlled comparison, but real semantic communication systems may operate on messages from different domains. Future experiments should test news text, conversational text, technical instructions, and short command-like sentences separately. This would reveal whether the learned semantic representation is general or whether it depends heavily on the training corpus distribution.
Error taxonomy construction: future work should classify decoding errors into categories such as synonym substitution, function-word error, named-entity error, tense error, word-order error, hallucinated phrase, and omission. Such an error taxonomy would make the analysis more interpretable than aggregate metrics alone. It would also help identify which errors should be handled by beam search, which by better channel modelling, and which by constrained language-model refinement.
In addition to these technical extensions, future work should also strengthen methodological breadth. Robustness claims can be reinforced through repeated runs with different random seeds and confidence-interval reporting for every metric. Such statistical reporting would make it easier to distinguish small true effects from run-to-run variance, especially when comparing close-performing variants at high SNR.

Another important direction is practical deployment profiling. While this report focuses on reconstruction quality, real deployed systems also care about latency, memory footprint, and throughput. A complete engineering evaluation should therefore include quality--cost curves, for example BLEU gain versus decoding time for greedy and beam settings, BLEU gain versus added parameters for different refinement schemes, and energy usage under repeated inference. This richer reporting would make the results directly actionable for system designers.

A longer-term research direction is to extend the semantic communication framework beyond text. Multimodal signals such as speech, images, and video carry significantly richer semantic content and are also more sensitive to channel-induced errors. Adapting the channel-aware mechanisms proposed in this project to multimodal transceivers would likely require more substantial architectural changes, but the underlying methodology, namely the hypothesis--failure--redesign loop with explicit negative-result reporting, can be reused directly.

\subsection{Final Remark}

This project shows that semantic communication transceiver design requires careful alignment between model architecture, channel condition, training objective, and decoding strategy. In the current experiments, the most reliable configuration is the Improved transceiver with Beam5 decoding and without external LLM refinement. At the same time, the mixed Rayleigh results, the semantic-loss ablation, and the FT-BERT refinement results show that several components still require further study. Future work should therefore focus on richer channel-state modelling, better-aligned semantic objectives, and more constrained refinement methods.

\FloatBarrier

\clearpage

\addcontentsline{toc}{section}{References}
\bibliographystyle{unsrt}
\bibliography{refs}
\clearpage

\appendix
\section{Appendices}
\subsection{Reproducible Command Flow}

This appendix collects the exact command-line invocations that generate all checkpoints and evaluation files referenced in the main text. All commands are executed from the project root on Windows PowerShell with the virtual environment activated. The expected output path of each training command is a checkpoint file under the \texttt{checkpoints/} tree, and the expected output path of each evaluation command is a variant-specific JSON file written next to the corresponding checkpoint.

\textbf{Environment setup.} The following commands create and activate a clean Python virtual environment and verify the CUDA availability that is required for reproducing all training runs on a single GPU workstation.
\begin{itemize}
    \item \texttt{python -m venv .venv}
    \item \texttt{.\textbackslash .venv\textbackslash Scripts\textbackslash Activate.ps1}
    \item \texttt{python -c}\\
    \texttt{\phantom{xx}"import torch; print(torch.cuda.is\_available())"}
\end{itemize}

\textbf{Training commands.} The following commands reproduce the five training runs reported in Tables~\ref{tab:awgn-main}, \ref{tab:ablation}, and \ref{tab:rayleigh}. The semantic-loss ablation is controlled by passing \texttt{--beta 0.0} and tagging the run with \texttt{--tag no\_sem}.
\begin{itemize}
    \item \texttt{python train.py}\\
    \texttt{\phantom{xx}--model baseline --channel awgn --epochs 80}
    \item \texttt{python train.py}\\
    \texttt{\phantom{xx}--model improved --channel awgn --epochs 80}
    \item \texttt{python train.py}\\
    \texttt{\phantom{xx}--model improved --channel awgn --epochs 80}\\
    \texttt{\phantom{xx}--beta 0.0 --tag no\_sem}
    \item \texttt{python train.py}\\
    \texttt{\phantom{xx}--model baseline --channel rayleigh --epochs 80}
    \item \texttt{python train.py}\\
    \texttt{\phantom{xx}--model improved --channel rayleigh --epochs 80}
\end{itemize}

\textbf{Evaluation commands.} The following evaluation commands reproduce the quantitative results reported in the main text. The \texttt{--beam\_width} flag switches decoding from greedy to beam search, and the \texttt{--use\_llm} flag together with \texttt{--llm\_model} enables the fine-tuned BERT refiner.
\begin{itemize}
    \item \texttt{python evaluate.py}\\
    \texttt{\phantom{xx}--checkpoint checkpoints/baseline\_awgn/best.pt}\\
    \texttt{\phantom{xx}--model baseline --channel awgn}
    \item \texttt{python evaluate.py}\\
    \texttt{\phantom{xx}--checkpoint checkpoints/improved\_awgn/best.pt}\\
    \texttt{\phantom{xx}--model improved --channel awgn}
    \item \texttt{python evaluate.py}\\
    \texttt{\phantom{xx}--checkpoint}\\
    \texttt{\phantom{xx}checkpoints/improved\_awgn\_no\_sem/best.pt}\\
    \texttt{\phantom{xx}--model improved --channel awgn}
    \item \texttt{python evaluate.py}\\
    \texttt{\phantom{xx}--checkpoint checkpoints/baseline\_rayleigh/best.pt}\\
    \texttt{\phantom{xx}--model baseline --channel rayleigh}
    \item \texttt{python evaluate.py}\\
    \texttt{\phantom{xx}--checkpoint checkpoints/improved\_rayleigh/best.pt}\\
    \texttt{\phantom{xx}--model improved --channel rayleigh}
    \item \texttt{python evaluate.py}\\
    \texttt{\phantom{xx}--checkpoint checkpoints/improved\_awgn/best.pt}\\
    \texttt{\phantom{xx}--model improved --channel awgn}\\
    \texttt{\phantom{xx}--beam\_width 5 --eval\_tag beam5}
    \item \texttt{python evaluate.py}\\
    \texttt{\phantom{xx}--checkpoint checkpoints/improved\_awgn/best.pt}\\
    \texttt{\phantom{xx}--model improved --channel awgn --use\_llm}\\
    \texttt{\phantom{xx}--llm\_model checkpoints/bert\_europarl}\\
    \texttt{\phantom{xx}--eval\_tag ftbert}
    \item \texttt{python evaluate.py}\\
    \texttt{\phantom{xx}--checkpoint checkpoints/improved\_awgn/best.pt}\\
    \texttt{\phantom{xx}--model improved --channel awgn --beam\_width 5}\\
    \texttt{\phantom{xx}--use\_llm}\\
    \texttt{\phantom{xx}--llm\_model checkpoints/bert\_europarl}\\
    \texttt{\phantom{xx}--eval\_tag beam5\_ftbert}
\end{itemize}

\subsection{Data Files Used for Quantitative Analysis}

All tables in the main text are computed directly from the following JSON files. Every file is generated by the corresponding evaluation command above, stored under the \texttt{checkpoints/} tree, and is amenable to version control.

\begin{itemize}
    \item \texttt{checkpoints/baseline\_awgn/eval\_awgn.json}
    \item \texttt{checkpoints/improved\_awgn/eval\_awgn.json}
    \item \texttt{checkpoints/improved\_awgn\_no\_sem/}\\
    \texttt{\phantom{xx}eval\_awgn.json}
    \item \texttt{checkpoints/baseline\_rayleigh/eval\_rayleigh.json}
    \item \texttt{checkpoints/improved\_rayleigh/eval\_rayleigh.json}
    \item \texttt{checkpoints/improved\_awgn/eval\_awgn\_beam5.json}
    \item \texttt{checkpoints/improved\_awgn/eval\_awgn\_ftbert.json}
    \item \texttt{checkpoints/improved\_awgn/}\\
    \texttt{\phantom{xx}eval\_awgn\_beam5\_ftbert.json}
\end{itemize}

\subsection{Notes on Figure and Table Preparation}

To satisfy the report template requirements, the following conventions are followed consistently throughout the document. First, every table and figure must be referenced in the main text before its appearance, so that the reader always encounters a pointer before the object itself. Second, captions are written to be informative and self-contained, so that a reader who only glances at the table can still interpret the main message. Third, axis labels and table column headers include units where relevant (for example \si{dB} for SNR), and every SNR sweep explicitly states the corresponding channel model so that readers cannot confuse AWGN and Rayleigh rows by accident.

For the final submission version, it is also recommended to include at least one figure that visualizes the overall experimental workflow (from data preprocessing through model training and evaluation), and one figure that compares key BLEU-4 trends across the major model variants across the full SNR range. Such visual summaries improve readability for examiners, compress the key message of multi-row tables, and make the evidence chain from design choice to quantitative outcome clearer. Future revisions of this thesis could add these figures without any change to the surrounding text, thanks to the figure-placement conventions described above.

\subsection{Reproducibility Check-list}

The following check-list summarises the minimum steps required to independently reproduce the quantitative claims of the report:
\begin{enumerate}
    \item Prepare the Europarl corpus and run the preprocessing script once to generate tokenized shards.
    \item Execute the five training commands listed above; each one writes a \texttt{best.pt} checkpoint under its own checkpoint folder.
    \item Execute the eight evaluation commands listed above; each one writes a JSON file that can be directly loaded to reproduce the corresponding row of the result tables.
    \item Cross-check the reported numbers in Tables 1--4 against the numerical fields of the JSON files, using the provided comparison helper script.
\end{enumerate}

Following this check-list end-to-end takes roughly two to three days on a single mid-range GPU, with the Rayleigh training runs accounting for the largest share of the total wall-clock time. All hyperparameters are documented either inside the training commands or inside the project configuration file, so that no further undocumented parameter is needed to reproduce the reported figures.

\FloatBarrier

\end{document}
